% =====================================================================
\chapter*{Instrucciones}
\phantomsection\addcontentsline{toc}{chapter}{Instrucciones}
\markboth{Instrucciones}{}
% =====================================================================

Este examen cierra la Fase~II del curso, \emph{Aprendizaje}, y evalúa los módulos~7 a~11. Comprueba que el estudiante sabe entrenar una red neuronal: plantear la pérdida, optimizarla, derivar sus gradientes a mano y con diferenciación automática, y mantener estable el entrenamiento.

\section*{Condiciones}
\begin{center}
\begin{tabular}{@{}ll@{}}
\toprule
Duración & 2~h 30~min \\
Material permitido & calculadora científica y una hoja A4 de fórmulas escrita a mano \\
Puntuación total & 100 puntos \\
Aprobado & 70 puntos (70\,\%) \\
Módulos evaluados & M07 Optimización, M08 Aprendizaje supervisado, \\
 & M09 Redes neuronales, M10 Diferenciación automática, \\
 & M11 Entrenamiento de redes profundas \\
\bottomrule
\end{tabular}
\end{center}

\section*{Estructura y puntuación}
\begin{itemize}[leftmargin=1.4em]
  \item \textbf{Parte A. Cuestiones} (30 puntos): 20 preguntas de opción múltiple, de 1,5 puntos cada una. Cada pregunta tiene una única respuesta correcta. Las respuestas erróneas no restan.
  \item \textbf{Parte B. Problemas} (70 puntos): 6 problemas de desarrollo, con la puntuación indicada en cada apartado. Deben escribirse todos los pasos: el planteamiento correcto puntúa aunque el resultado numérico sea erróneo, y un error aritmético arrastrado de forma coherente solo se penaliza en el apartado en que se comete.
  \item El \textbf{proyecto de fase} se entrega aparte, en un plazo de dos semanas. En la nota de la fase, el examen pesa un 40\,\% y el proyecto un 30\,\%.
\end{itemize}

\section*{Convenios}
\begin{itemize}[leftmargin=1.4em]
  \item Salvo indicación en contrario, se redondea a cuatro cifras decimales y se usan logaritmos naturales (las pérdidas se expresan en nats).
  \item \textbf{Vectores fila}, como en NumPy y PyTorch: $z^{(1)}=xW_1+b_1$, $h=\varphi(z^{(1)})$, $z^{(2)}=hW_2+b_2$, $p=\softmax(z^{(2)})$ y $L=-\ln p_y$. Con un lote, $x$ pasa a ser la matriz $X$ de tamaño $B\times n$ y la softmax se aplica por filas.
  \item Los \emph{deltas} son $\delta^{(l)}=\partial L/\partial z^{(l)}$; con un lote, $\Delta^{(l)}$ apila los deltas de los ejemplos por filas.
  \item En un grafo computacional, $\bar v=\partial L/\partial v$ es el \emph{adjunto} del nodo $v$ (M10, §3.1).
  \item Las clases se numeran $0,1,2,\dots$ y $e_k$ es el vector canónico con un~1 en la posición~$k$.
  \item $\sigma(t)=1/(1+e^{-t})$ es la sigmoide y $\sigma'(t)=\sigma(t)\bigl(1-\sigma(t)\bigr)$.
\end{itemize}

\section*{Procedencia de los problemas}
Cuatro de los seis problemas de la parte~B (los problemas 1 a~4) están adaptados de exámenes y hojas de problemas de cursos universitarios: dos exámenes parciales y un examen final de CS189 de Berkeley y la primera tarea de CS224N de Stanford; un apartado del problema~5 procede también de esta última. Cada problema indica su \emph{Fuente}. Los enunciados son de redacción propia y los datos se han elegido para que puedan resolverse a mano. Las soluciones completas están en el \cref{ap:soluciones} y la tabla de autoevaluación en el \cref{ap:autoevaluacion}.
